\documentclass{article}
\usepackage[T1]{fontenc}
\usepackage{lmodern}
\usepackage{graphicx}
\usepackage{amsmath,amssymb}
\usepackage[a4paper,margin=2cm]{geometry}
\usepackage[absolute,overlay]{textpos}
\setlength{\TPHorizModule}{1mm}
\setlength{\TPVertModule}{1mm}
\usepackage[indonesian]{babel}
\usepackage{hyperref}
\setlength{\emergencystretch}{2em}
% unit: Halaman 52

\begin{document}

% === Halaman 52 (llm) ===

\section*{Menghitung Ukuran Pesan yang dapat Disembunyikan}

\begin{itemize}
    \item Ukuran pesan yang akan disembunyikan bergantung pada ukuran \textit{cover-object}.
    \item Misalkan pada citra \textit{grayscale} (1 byte/pixel) $256 \times 256$ \textit{pixel} :
    \begin{itemize}
        \item jumlah \textit{pixel} = jumlah byte = $256 \times 256 = 65536$
        \item setiap byte dapat menyembunyikan 1 bit pesan di LSB-nya
        \item jadi ukuran maksimal pesan = $65536$ bit = $8192$ \textit{byte} = 8 KB
    \end{itemize}
    \item Pada citra berwarna 24-bit berukuran $256 \times 256$ \textit{pixel}:
    \begin{itemize}
        \item jumlah \textit{pixel} $256 \times 256 = 65536$
        \item setiap pixel = 3 byte, berarti ada $65536 \times 3 = 196608$ \textit{byte}.
        \item setiap \textit{byte} dapat menyembunyikan 1 bit pesan
        \item jadi ukuran maksimal pesan = $196608$ bit = $24576$ byte = 24KB
    \end{itemize}
\end{itemize}

\end{document}
